\documentclass[10pt]{article}

\usepackage[utf8]{inputenc}

\usepackage[T1]{fontenc}

\usepackage{amsmath}

\usepackage{amsfonts}

\usepackage{amssymb}

\usepackage[version=4]{mhchem}

\usepackage{stmaryrd}

\usepackage{graphicx}

\usepackage[export]{adjustbox}

\graphicspath{ {./images/} }

\usepackage{mathrsfs}



\DeclareUnicodeCharacter{0131}{$\imath$}



\begin{document}

\begin{enumerate}

  \setcounter{enumi}{3}

  \item Неотрицательные целые числа, упорядоченные по делимости: $a \leqslant b$, если $b=a c$ для нек-рого $c$, где sup $\{a$, $b\}$ - наименьітее общее кратное $a$ и $b, \mathrm{a} \inf \{a, b\}-$ наибольший общий делитель $a$ и $b$.



  \item Действительные функции, определенные на отрезке $[0,1]$ и упорядоченные условием: $f \leqslant g$, если $(t) \leqslant g(t)$ для всех $t \in[0,1]$, где



\end{enumerate}



\[

\sup \{f, g\}=u

\]



причем



\[

u(t)=\max \{f(t), g(t)\}

\]



a



\[

\inf \{f, g\}=v

\]



причем



\[

v(t)=\min \{f(t), \quad g(t)\} .

\]



Пусть $M$ - решетка. $M$ становится универсальной алгеброй с двумя бинарными операциями, если определить



\[

\begin{gathered}

a+b=\sup \{a, b\}, \\

a \cdot b=\inf \{a, b\}

\end{gathered}

\]



(вместо + и - часто употребляются символы $U$ и $\cap$ или $\checkmark$ и $\wedge$ ). Эта универсальная алгебра удовлетворяет следующим тождественным соотношениям:



(1) $a+a=a$



$\left(1^{\prime}\right) a \cdot a=a$



(2) $a+b=b+a$;



$\left(2^{\prime}\right) a \cdot b=b \cdot a$;



(3) $(a+b)+c=a+(b+c) ;\left(3^{\prime}\right)(a \cdot b) \cdot c=a \cdot(b \cdot c)$;



(4) $a(a+b)=a$;



(4') $a+a \cdot b=a$.



Наоборот, если $M$ - множество с двумя бинарными операциями, обладающими перечислешными выше свойствами (1) - (4), $\left(1^{\prime}\right)-\left(4^{\prime}\right)$, то на $M$ можно задать порядок $\leqslant$, полагая $a \leqslant b$, если $a+b=b$ (при этом окажется, что $a \leqslant b$ тогда и только тогда, когда $a \cdot b=a)$. Возникающее частично упорядоченное множество будет Р., причем



\[

\sup \{a, b\}=a+b, \mathrm{a} \inf \{a, b\}=a \cdot b .

\]



Таким образом, Р. можно определить как универсальную алгебру, описываемую тождествами (1)-(4), (1') $\left(4^{\prime}\right)$, то есть Р. образует универсальных алгебр многообразие.



Если частично упорядоченное множество рассматривать как малую категорию, то оно оказывается Р. в том и только в том случае, когда для любых двух объектов этой категории существует их произведение и копроизведение.



Если $P$ и $P^{\prime}$ - решетки и $f$ - изоморфизм этих частично упорядоченных множеств, то $f$ является также изоморфизмом соответствующих универсальных алгебр, т. е.



\[

f(x+y)=f(x)+f(y) \text { и } f(x y)=f(x) \cdot f(y)

\]



для любых $x, y \in P$. Однако произвольное изотонное отображение решетки $P$ в решетку $P^{\prime}$ не обязано быть гомоморфизмом этих $P$., рассматриваемых как универсальные алгебры. Так, для любого $a \in P$ отображения $f(x)=x+a$ и $g(x)=x \cdot a$ - изотонные отображения решетки $P$ в себя, являющиеся гомоморфизмами лишь в том и только в том случае, когда $P$ - дистрибутивная решетка. Впрочем, первое из этих отображений является гомоморфизмом полурешетки $P$ с операцией + , а второе-гомоморфизмом полурешетки $P$ с операцией · Совокупность всех Р. образует категорию, если морфизмами считать гомоморфизмы.



Антитомоморфизм ре шіетки $P$ в решетку $P^{\prime}$ есть такое отображение $f: P \rightarrow P^{\prime}$, что



\[

f(x+y)=f(x) \cdot f(y) \text { и } f(x \cdot y)=f(x)+f(y)

\]



для любых $x, y \in P$. Последовательное выполнение двух антигомоморфизмов является гомоморфизмом. Частично упорядоченное множество, антиизоморфное Р., есть Р. Под к о о р д и н а т и з а ц и е й Р. понимают нахождение алгебрачческой системы (чаще, универсальной алгебры) такой, что данная Р. изоморфна Р. подсистем, Р. конгруэнций или какой-либо другой Р., связанной с этой алгебраич. системой или универсальной алгеброй. Произвольная Р. с 0 и 1 координатизируется частично упорядоченной полугруппой ее резидуальных отображений в себя, оказываясь изоморфной Р. правых аннуляторов этой полугруппы. Сама полугруппа является б э р о в с к о й, т. е. как правый, так и левый аннулятор каждого из ее элементов порождается идемпотентом.



Наиболее важныө результаты получены для Р., подчиненных тем или иным дополнительным ограничениям (см. Алгебрачческая решетка, А томнал решетка, Брауэра решетка, Векторная решетка, Дедекиндова решетка, Дистрибутивнал решетка, Мультипликативнал решетка, Ортамодулярная решетка, Полная решетка, Свободная решетка, Решетпа с дополнениями, Булева алгебра). По отдельным вопросам теории Р. см. Идеал, Фильтр, Пополнение сеченияли. Особую роль играют алгебраич. образования, являющиеся в то же время Р. (см. Структурно упорядоченная группа). Наибольшее число приложений теории Р. связано с булевыми алгебрами. Другие классы Р. использовались в квантовой механике и физике.



Появление понятия «Р.» относится к сер. 19 в. и связано с тем, что многие факты, касающиеся множества идеалов кольца и множества нормальных подгрупп групшы, выглядят аналогично и могут быть доказаны в рамках теории дедекиндовых Р. Как самостоятельное направление алгебры теория $P$. сформировалась в 30-х гг. 20 в.



Лит.: [1] Б и р к г о ф Г., Теория решеток, пер. с англ., м., 1983; [2] $\Gamma$ p e т ц е p Саратов, 1970; [4] С к о р н я к о в Л. А., Дедекиндовы струкСаратов, 1970; [4] С к о р н я к о в Л. А., Дедекиндовы струкж е, Элементы теории структур, 2 и зд., М., 1982 ; [6] Йтоги науки. Алгебра. 1964, М., 1966, с. 237-74, [7]' Итоги науки. Алгебра, топология, геометрия. 1968, М., 1970, с. 101-54; [8] УпоB l y t h T., J a n o w i t z M., Resıduation theory, N. Y., 1972.



\includegraphics[max width=\textwidth, center]{2023_07_08_d39ed41f86a0480194f6g-1}

а л г е б р ы $A$ - частично упорядоченное (отношениөм төорөтико-множественного включения) множество $\mathrm{Sub} A$ всех подалгебр алгебры $A$. Для произвольных $X, Y \in \operatorname{Sub} A$ их супрөмумом будет подалгебра, порожденная $X$ и $Y$, а их инфинумом - пересечение $X \cap Y$. Пересечение подалгебр может быть пустым, поәтому для нек-рых типов алгебр (напр., для полугрупп и репшеток) к числу подалгебр относят и пустое множество. Для любой алгебры $A$ P. п. Sub $A$ являөтся алгебраической и обратно, для любой алгөбраической решетки $L$ существует алгебра $A$ такая, что $L \cong \operatorname{Sub} A$ (т ө о р е м а Би р к гоф а - Фри н к а). Любая решетка вложима в решетку $\operatorname{Sub} A$ для нек-рой группы $A$.



P. п. Sub $A$ - одна из основных производных структур, сопоставляемых алгебре $A$ (наряду с такими структурами, как группа автоморфизмов, полугруппа эндоморфизмов, решетка конгруэнций и т. п.). Проблематика, посвященная изучению связей между алгебрами и их Р. п. делится на такие аспекты: решеточные изоморфизмы, решеточные характеристики тех или иных классов алгебр, исследование алгебр с различными ограничениями на Р. п. Алгебры $A$ и $B$ наз. р е ш е т о ч н о и з о м о р ф н м м и, если $\operatorname{Sub} A \cong \operatorname{Sub} B$; всякий изоморфизм $\operatorname{Sub} A$ на $\operatorname{Sub} B$ наз. р е іI е т о ч ны м и зомо р ф имом (или проекти рован ием) $A$ на $B$. Изоморфные алгебры решеточно изоморфны, обратное же далеко не обязательно. Говорят, что алгебра $A$ р е іI е т о ч о о п р е д е л я е т с я (в данном классе $\mathscr{\kappa}$ ), если для любой однотипной с ней алгебры $B$ (из $\mathscr{K})$ из $\mathrm{Sub} B \cong \operatorname{Sub} A$ следует $B \cong A$. В некото-





\end{document}